\begin{thebibliography}{99}
\bibitem{r1} N. Aletras, D. Tsarapatsanis, D. Preo\c{t}iuc-Pietro, and V. Lampos, ``Predicting judicial decisions of the European Court of Human Rights: A natural language processing perspective,'' \emph{PeerJ Computer Science}, vol. 2, e93, 2016.
\bibitem{r2} M. Medvedeva, M. Vols, and M. Wieling, ``Using machine learning to predict decisions of the European Court of Human Rights,'' \emph{Artificial Intelligence and Law}, vol. 28, pp. 237--266, 2020.
\bibitem{r3} I. Chalkidis, I. Androutsopoulos, and N. Aletras, ``Neural legal judgment prediction in English,'' in \emph{Proc. 57th Annual Meeting of the Association for Computational Linguistics (ACL)}, 2019, pp. 4317--4323.
\bibitem{r4} D. M. Katz, M. J. Bommarito II, and J. Blackman, ``A general approach for predicting the behavior of the Supreme Court of the United States,'' \emph{PLOS ONE}, vol. 12, no. 4, e0174698, 2017.
\bibitem{r5} H. Zhong, Z. Guo, C. Tu, C. Xiao, Z. Liu, and M. Sun, ``Legal judgment prediction via topological learning,'' in \emph{Proc. Conference on Empirical Methods in Natural Language Processing (EMNLP)}, 2018, pp. 3540--3549.
\bibitem{r6} C. Xiao, H. Zhong, Z. Guo, C. Tu, Z. Liu, M. Sun, et al., ``CAIL2018: A large-scale legal dataset for judgment prediction,'' arXiv:1807.02478, 2018.
\bibitem{r7} J. Niklaus, I. Chalkidis, and M. St\"urmer, ``Swiss-Judgment-Prediction: A multilingual legal judgment prediction benchmark,'' in \emph{Proc. Natural Legal Language Processing Workshop (NLLP)}, 2021.
\bibitem{r8} H. Zhong, C. Xiao, C. Tu, T. Zhang, Z. Liu, and M. Sun, ``How does NLP benefit legal system: A summary of legal artificial intelligence,'' in \emph{Proc. 58th Annual Meeting of the ACL}, 2020, pp. 5218--5230.
\bibitem{r9} Y. Feng, C. Li, and V. Ng, ``Legal judgment prediction: A survey of the state of the art,'' in \emph{Proc. 31st International Joint Conference on Artificial Intelligence (IJCAI), Survey Track}, 2022, pp. 5461--5469.
\bibitem{r10} V. Malik, R. Sanjay, S. K. Nigam, K. Ghosh, S. K. Guha, A. Bhattacharya, and A. Modi, ``ILDC for CJPE: Indian legal documents corpus for court judgment prediction and explanation,'' in \emph{Proc. ACL-IJCNLP}, 2021, pp. 4046--4062.
\bibitem{r11} A. Joshi, S. Paul, A. Sharma, P. Goyal, S. Ghosh, and A. Modi, ``IL-TUR: Benchmark for Indian legal text understanding and reasoning,'' in \emph{Proc. 62nd Annual Meeting of the ACL}, 2024.
\bibitem{r12} S. Paul, A. Mandal, P. Goyal, and S. Ghosh, ``Pre-trained language models for the legal domain: A case study on Indian law,'' in \emph{Proc. 19th International Conference on Artificial Intelligence and Law (ICAIL)}, 2023.
\bibitem{r13} E. J. Hu, Y. Shen, P. Wallis, Z. Allen-Zhu, Y. Li, S. Wang, L. Wang, and W. Chen, ``LoRA: Low-rank adaptation of large language models,'' in \emph{Proc. International Conference on Learning Representations (ICLR)}, 2022.
\bibitem{r14} T. Dettmers, A. Pagnoni, A. Holtzman, and L. Zettlemoyer, ``QLoRA: Efficient finetuning of quantized LLMs,'' in \emph{Advances in Neural Information Processing Systems (NeurIPS)}, 2023.
\bibitem{r15} N. Xu, P. Wang, L. Chen, L. Pan, X. Wang, and J. Zhao, ``Distinguish confusing law articles for legal judgment prediction,'' in \emph{Proc. 58th Annual Meeting of the ACL}, 2020, pp. 3086--3095.
\bibitem{r16} S. Paul, P. Goyal, and S. Ghosh, ``LeSICiN: A heterogeneous graph-based approach for automatic legal statute identification from Indian legal documents,'' in \emph{Proc. AAAI Conference on Artificial Intelligence}, vol. 36, 2022.
\bibitem{r17} P. Bhattacharya, K. Ghosh, A. Pal, and S. Ghosh, ``Hier-SPCNet: A legal statute hierarchy-based heterogeneous network for computing legal case document similarity,'' in \emph{Proc. 43rd International ACM SIGIR Conference}, 2020.
\bibitem{r18} Q. Dong and S. Niu, ``Legal judgment prediction via relational learning,'' in \emph{Proc. 44th International ACM SIGIR Conference on Research and Development in Information Retrieval}, 2021.
\bibitem{r19} M. Yasunaga, H. Ren, A. Bosselut, P. Liang, and J. Leskovec, ``QA-GNN: Reasoning with language models and knowledge graphs for question answering,'' in \emph{Proc. NAACL-HLT}, 2021, pp. 535--546.
\bibitem{r20} X. Zhang, A. Bosselut, M. Yasunaga, H. Ren, P. Liang, C. D. Manning, and J. Leskovec, ``GreaseLM: Graph reasoning enhanced language models for question answering,'' in \emph{Proc. International Conference on Learning Representations (ICLR)}, 2022.
\bibitem{r21} B. Perozzi, B. Fatemi, D. Zelle, A. Tsitsulin, M. Kazemi, R. Al-Rfou, and J. Halcrow, ``Let your graph do the talking: Encoding structured data for LLMs,'' arXiv:2402.05862, 2024.
\bibitem{r22} Y. Tian, H. Song, Z. Wang, H. Wang, Z. Hu, F. Wang, N. V. Chawla, and P. Xu, ``Graph neural prompting with large language models,'' in \emph{Proc. AAAI Conference on Artificial Intelligence}, vol. 38, 2024.
\bibitem{r23} X. He, Y. Tian, Y. Sun, N. V. Chawla, T. Laurent, Y. LeCun, X. Bresson, and B. Hooi, ``G-Retriever: Retrieval-augmented generation for textual graph understanding and question answering,'' in \emph{Advances in Neural Information Processing Systems (NeurIPS)}, 2024.
\bibitem{r24} M. Medvedeva and P. McBride, ``Legal judgment prediction: If you are going to do it, do it right,'' in \emph{Proc. Natural Legal Language Processing Workshop (NLLP)}, 2023.
\bibitem{r25} A. S{\o}gaard, S. Ebert, J. Bastings, and K. Filippova, ``We need to talk about random splits,'' in \emph{Proc. 16th Conference of the European Chapter of the ACL (EACL)}, 2021, pp. 1823--1832.
\bibitem{r26} A. Lazaridou et al., ``Mind the gap: Assessing temporal generalization in neural language models,'' in \emph{Advances in Neural Information Processing Systems (NeurIPS)}, 2021.
\bibitem{r27} S. Kapoor and A. Narayanan, ``Leakage and the reproducibility crisis in machine-learning-based science,'' \emph{Patterns}, vol. 4, no. 9, 100804, 2023.
\bibitem{r28} R. Dror, G. Baumer, S. Shlomov, and R. Reichart, ``The hitchhiker's guide to testing statistical significance in natural language processing,'' in \emph{Proc. 56th Annual Meeting of the ACL}, 2018, pp. 1383--1392.
\bibitem{r32} M. Schlichtkrull, T. N. Kipf, P. Bloem, R. van den Berg, I. Titov, and M. Welling, ``Modeling relational data with graph convolutional networks,'' in \emph{Proc. Extended Semantic Web Conference (ESWC)}, 2018, pp. 593--607.
\bibitem{r31} H. Small, ``Co-citation in the scientific literature: A new measure of the relationship between two documents,'' \emph{Journal of the American Society for Information Science}, vol. 24, no. 4, pp. 265--269, 1973.
\bibitem{r36} L. Yue, Q. Liu, B. Jin, H. Wu, K. Zhang, Y. An, M. Cheng, B. Yin, and D. Wu, ``NeurJudge: A circumstance-aware neural framework for legal judgment prediction,'' in \emph{Proc. 44th International ACM SIGIR Conference on Research and Development in Information Retrieval}, 2021.
\bibitem{r37} Y. Feng, C. Li, and V. Ng, ``Legal judgment prediction via event extraction with constraints,'' in \emph{Proc. 60th Annual Meeting of the ACL}, 2022, pp. 648--664.
\bibitem{r38} I. Chalkidis, A. Jana, D. Hartung, M. Bommarito, I. Androutsopoulos, D. M. Katz, and N. Aletras, ``LexGLUE: A benchmark dataset for legal language understanding in English,'' in \emph{Proc. 60th Annual Meeting of the ACL}, 2022.
\bibitem{r39} J. Niklaus, V. Matoshi, P. Rani, A. Galassi, M. St\"urmer, and I. Chalkidis, ``LEXTREME: A multi-lingual and multi-task benchmark for the legal domain,'' in \emph{Findings of the Association for Computational Linguistics: EMNLP 2023}, 2023.
\bibitem{r40} I. Chalkidis, M. Fergadiotis, D. Tsarapatsanis, N. Aletras, I. Androutsopoulos, and P. Malakasiotis, ``Paragraph-level rationale extraction through regularization: A case study on European Court of Human Rights cases,'' in \emph{Proc. NAACL-HLT}, 2021, pp. 226--241.
\bibitem{r41} P. Kalamkar, A. Tiwari, A. Agarwal, S. Karn, S. Gupta, V. Raghavan, and A. Modi, ``Corpus for automatic structuring of legal documents,'' in \emph{Proc. 13th Language Resources and Evaluation Conference (LREC)}, 2022.
\bibitem{r42} S. Vats, A. Zope, S. De, A. Sharma, U. Bhattacharya, S. K. Nigam, S. Guha, K. Rudra, and K. Ghosh, ``LLMs -- the good, the bad or the indispensable?: A use case on legal statute prediction and legal judgment prediction on Indian court cases,'' in \emph{Findings of the Association for Computational Linguistics: EMNLP 2023}, 2023.
\bibitem{r43} S. K. Nigam, A. Sharma, D. Khanna, N. Shallum, K. Ghosh, and A. Bhattacharya, ``Legal judgment reimagined: PredEx and the rise of intelligent AI interpretation in Indian courts,'' in \emph{Findings of the Association for Computational Linguistics: ACL 2024}, 2024.
\bibitem{r44} I. Chalkidis, M. Fergadiotis, P. Malakasiotis, N. Aletras, and I. Androutsopoulos, ``LEGAL-BERT: The Muppets straight out of law school,'' in \emph{Findings of the Association for Computational Linguistics: EMNLP 2020}, 2020, pp. 2898--2904.
\bibitem{r45} P. Henderson, M. S. Krass, L. Zheng, N. Guha, C. D. Manning, D. Jurafsky, and D. E. Ho, ``Pile of Law: Learning responsible data filtering from the law and a 256GB open-source legal dataset,'' in \emph{Advances in Neural Information Processing Systems (NeurIPS), Datasets and Benchmarks Track}, 2022.
\bibitem{r46} S. Gururangan, A. Marasovi\'{c}, S. Swayamdipta, K. Lo, I. Beltagy, D. Downey, and N. A. Smith, ``Don't stop pretraining: Adapt language models to domains and tasks,'' in \emph{Proc. 58th Annual Meeting of the ACL}, 2020, pp. 8342--8360.
\bibitem{r47} P. Colombo, T. P. Pires, M. Boudiaf, D. Culver, R. Melo, C. Corro, A. F. T. Martins, F. Esposito, V. L. Raposo, S. Morgado, and M. Desa, ``SaulLM-7B: A pioneering large language model for law,'' arXiv:2403.03883, 2024.
\bibitem{r48} N. Guha, J. Nyarko, D. E. Ho, C. R\'{e}, et al., ``LegalBench: A collaboratively built benchmark for measuring legal reasoning in large language models,'' in \emph{Advances in Neural Information Processing Systems (NeurIPS), Datasets and Benchmarks Track}, 2023.
\bibitem{r49} R. Shui, Y. Cao, X. Wang, and T.-S. Chua, ``A comprehensive evaluation of large language models on legal judgment prediction,'' in \emph{Findings of the Association for Computational Linguistics: EMNLP 2023}, 2023.
\bibitem{r50} Y. Wu, S. Zhou, Y. Liu, W. Lu, X. Liu, Y. Zhang, C. Sun, F. Wu, and K. Kuang, ``Precedent-enhanced legal judgment prediction with LLM and domain-model collaboration,'' in \emph{Proc. Conference on Empirical Methods in Natural Language Processing (EMNLP)}, 2023.
\bibitem{r51} M. Dahl, V. Magesh, M. Suzgun, and D. E. Ho, ``Large legal fictions: Profiling legal hallucinations in large language models,'' \emph{Journal of Legal Analysis}, vol. 16, no. 1, 2024.
\bibitem{r52} J. H. Fowler, T. R. Johnson, J. F. Spriggs II, S. Jeon, and P. J. Wahlbeck, ``Network analysis and the law: Measuring the legal importance of precedents at the U.S. Supreme Court,'' \emph{Political Analysis}, vol. 15, no. 3, pp. 324--346, 2007.
\bibitem{r53} A. Hogan, E. Blomqvist, M. Cochez, C. d'Amato, G. de Melo, C. Gutierrez, et al., ``Knowledge graphs,'' \emph{ACM Computing Surveys}, vol. 54, no. 4, pp. 1--37, 2021.
\bibitem{r54} T. N. Kipf and M. Welling, ``Semi-supervised classification with graph convolutional networks,'' in \emph{Proc. International Conference on Learning Representations (ICLR)}, 2017.
\bibitem{r55} M. X. Goemans and D. P. Williamson, ``A general approximation technique for constrained forest problems,'' \emph{SIAM Journal on Computing}, vol. 24, no. 2, pp. 296--317, 1995.
\bibitem{r56} P. Lewis, E. Perez, A. Piktus, F. Petroni, V. Karpukhin, N. Goyal, H. K\"uttler, M. Lewis, W. Yih, T. Rockt\"aschel, S. Riedel, and D. Kiela, ``Retrieval-augmented generation for knowledge-intensive NLP tasks,'' in \emph{Advances in Neural Information Processing Systems (NeurIPS)}, 2020.
\bibitem{r57} D. Edge, H. Trinh, N. Cheng, J. Bradley, A. Chao, A. Mody, S. Truitt, and J. Larson, ``From local to global: A graph RAG approach to query-focused summarization,'' arXiv:2404.16130, 2024.
\bibitem{r58} S. Pan, L. Luo, Y. Wang, C. Chen, J. Wang, and X. Wu, ``Unifying large language models and knowledge graphs: A roadmap,'' \emph{IEEE Transactions on Knowledge and Data Engineering}, vol. 36, no. 7, 2024.
\bibitem{r33} X. L. Li and P. Liang, ``Prefix-tuning: Optimizing continuous prompts for generation,'' in \emph{Proc. ACL-IJCNLP}, 2021, pp. 4582--4597.
\bibitem{r59} A. Yang et al. (Qwen Team), ``Qwen3 technical report,'' arXiv:2505.09388, 2025.
\bibitem{r34} B. Efron and R. J. Tibshirani, \emph{An Introduction to the Bootstrap}. New York: Chapman \& Hall/CRC, 1993.
\bibitem{r35} Q. McNemar, ``Note on the sampling error of the difference between correlated proportions or percentages,'' \emph{Psychometrika}, vol. 12, no. 2, pp. 153--157, 1947.
\bibitem{r60} Y. Cui, M. Jia, T.-Y. Lin, Y. Song, and S. Belongie, ``Class-balanced loss based on effective number of samples,'' in \emph{Proc. IEEE/CVF Conference on Computer Vision and Pattern Recognition (CVPR)}, 2019, pp. 9268--9277.
\bibitem{r29} ``Indian-law-supreme-court-judgements-2016,'' Hugging Face dataset, Shreyasrao/Indian-law-supreme-court-judgements-2016, accessed Oct. 2026. [Online]. Available: \url{https://huggingface.co/datasets/Shreyasrao/Indian-law-supreme-court-judgements-2016}
\bibitem{r61} Y. Du et al., ``PP-OCR: A practical ultra lightweight OCR system,'' arXiv:2009.09941, 2020.
\bibitem{r30} ``Indian Supreme Court Judgments,'' Registry of Open Data on AWS, \texttt{s3://indian-supreme-court-judgments}, licence CC BY 4.0, accessed Oct. 11, 2026.
\bibitem{r62} CivicTech India, Indian law JSON datasets (Indian Penal Code, Code of Criminal Procedure, Indian Evidence Act, Constitution of India), GitHub repositories, commit-pinned, accessed October 2026.
\bibitem{r63} S. Robertson and H. Zaragoza, ``The probabilistic relevance framework: BM25 and beyond,'' \emph{Foundations and Trends in Information Retrieval}, vol. 3, no. 4, pp. 333--389, 2009.
\bibitem{r64} T. Wolf, L. Debut, V. Sanh, J. Chaumond, C. Delangue, A. Moi, P. Cistac, T. Rault, R. Louf, M. Funtowicz, et al., ``Transformers: State-of-the-art natural language processing,'' in \emph{Proc. EMNLP: System Demonstrations}, 2020, pp. 38--45.
\end{thebibliography}
